\chapter{相关理论与关键技术}
\label{chp:related_tech}

\section{图推理任务基础}

本节从图表示、典型图算法与任务特性出发，说明图推理问题与一般文本推理任务在输入形式、约束结构和求解目标上的差异，为后续方法设计提供基础定义。

\subsection{图与图算法}

设图记为 $G=(V,E)$，其中 $V$ 表示节点集合，$E$ 表示边集合。根据边是否具有方向和权重，还可区分为有向图、无向图、带权图、多重图以及动态图等形式\cite{diestel2017graph}。在众多图计算任务中，最短路径、连通分量、拓扑排序、最大流、匹配以及团搜索等问题具有较高代表性。

传统图算法在上述任务上已经形成了较为成熟的理论体系。例如，最短路径计算可采用 Dijkstra 或 Bellman-Ford，DFS 与 BFS 可用于遍历和连通性判断，Ford-Fulkerson 与 Edmonds-Karp 可用于流优化\cite{cormen2009clrs}，回溯搜索与动态规划则常用于处理旅行商等组合优化问题\cite{korte2018combinatorial}。从更广义的图学习视角看，DeepWalk\cite{perozzi2014deepwalk}和 node2vec\cite{grover2016node2vec}推动了基于随机游走的节点表征学习；GCN\cite{kipf2017gcn}、GraphSAGE\cite{hamilton2017graphsage}、GAT\cite{velickovic2018gat}和 GIN\cite{xu2019gin}发展了不同形式的消息传递网络；GraphGPS\cite{rampasek2022graphgps}则进一步融合图神经网络与 Transformer 架构。SNAP\cite{leskovec2016snap}、OGB\cite{hu2020ogb}、Benchmarking Graph Neural Networks\cite{dwivedi2023benchmarking}和 GraphGym\cite{you2020graphgym}分别为图数据组织、开放基准、模型比较和设计空间评测提供了基础。这些算法与基准共同说明，图任务具有清晰的结构对象、明确的计算过程和可复核的评测方式，适合与外部程序执行和结果验证机制结合。

\subsection{图推理任务的特点}

与一般文本推理相比，图推理任务同时具有强结构依赖、强约束驱动和强程序化特征。许多问题的答案并不来自局部语义片段，而是来自节点与边在全局拓扑中的相对位置；同时，大量任务还要求结果必须满足路径连通、回路闭合、容量限制、覆盖约束或最优性等严格条件。这意味着图推理不能仅满足叙述层面的自洽，还必须满足可计算正确性与约束一致性。

另一方面，图任务又具有鲜明的可执行性。最短路径、连通分量、拓扑排序、最大流、匹配和最小生成树等问题都存在成熟算法，很多组合优化问题也存在清晰的搜索框架。因此，相比把所有求解压力压在语言模型一次性生成答案上，更合理的思路是让模型承担任务理解、规划组织与知识调用，再将精确计算交给外部程序环境完成。

然而，图任务并不天然适合当前主流大语言模型的输入方式。语言模型擅长处理线性序列，但图结构本质上是离散关系对象的集合，需要在文本中额外编码节点、边、权重和约束。当图规模增大或约束变复杂时，文本化表示很容易失去局部相邻关系与全局拓扑的直观性，这也是图推理性能不稳定的重要原因。

\begin{figure}[htbp]
\centering
\includegraphics[width=0.96\textwidth]{figures/paper/graph_reasoning_challenge.pdf}
\caption{图推理任务的困难传递链}
\label{fig:graph_reasoning_challenge}
\end{figure}

图~\ref{fig:graph_reasoning_challenge} 概括了图推理任务中错误从输入端向输出端传递的过程。可以看到，如果输入阶段没有充分保留图结构信息，后续的任务识别、算法选择和程序实现都会受到影响，并可能进一步放大初始偏差。因此，中间规划表示的作用不仅在于帮助模型展开推理，更在于为任务语义、图结构约束和程序生成之间建立一个可检查的过渡层，从而减少错误在不同阶段之间的传递。
\subsection{图表示序列化}

图结构如何表示，是大语言模型处理图推理任务时首先需要面对的问题。现有任务中常见的表示方式包括边列表、邻接表、自然语言描述以及类似 JSON 的半结构化格式。边列表能够直接列出节点之间的连接关系，适合表达图中的具体边；但当边同时带有方向、权重、时间戳等属性时，输入长度和阅读难度都会增加。邻接表更突出节点的局部邻域信息，便于观察相邻节点关系，但对全局路径约束、连通关系或组合优化目标的表达不够直接。自然语言描述与用户提问形式更接近，容易承载任务背景和需求说明，但也可能把图结构、背景信息和求解目标混在一起，增加模型解析难度。半结构化表示更便于程序读取和后续解析，不过相较于自然语言，它对任务语义和隐含条件的表达能力相对有限。

更关键的是，序列化并非仅是格式转换过程，它会直接改变模型的注意力分配方式。在线性上下文中，相距较远的边信息即便在原图中高度相关，也可能因为跨越多个文本片段而难以被模型联合建模。对人类而言，图像化拓扑结构通常可以一眼看出关键连通关系；但对语言模型而言，同样的结构必须通过多轮文本匹配才能恢复。正因为如此，很多图推理基准测试集虽然规模不算大，却依然对模型构成明显挑战。

\begin{table}[htbp]
\centering
\small
\caption{常见图表示方式及其适用特点}
\label{tab:graph_representation}
\renewcommand{\arraystretch}{1.15}
\begin{tabularx}{\textwidth}{
>{\raggedright\arraybackslash}p{2.6cm}
>{\raggedright\arraybackslash}X
>{\raggedright\arraybackslash}X}
\toprule
表示方式 & 优点 & 局限性 \\
\midrule
边列表表示 & 关系清晰，便于精确列边 & 全局拓扑不直观，长序列输入负担大 \\
邻接表表示 & 突出局部邻接关系 & 全局路径和复杂约束表达不足 \\
自然语言描述 & 用户友好，便于嵌入任务背景 & 语义噪声较多，算法意图不集中 \\
半结构化表示 & 易于程序解析和转换 & 自然性较弱，可读性相对不足 \\
\bottomrule
\end{tabularx}
\end{table}

表\ref{tab:graph_representation}说明，不同表示方式各有优势，但没有一种方式可以同时兼顾自然表达、结构完整和实现友好。因此，很多高性能图推理系统并不会只依赖输入表示本身，而是会在读取输入后先构造结构化中间表示，再进入规划与程序生成阶段。

\section{大语言模型推理增强技术}

本节围绕分步推理、结构化搜索与推理过程控制等方向，归纳大语言模型在复杂任务中的能力增强路径，并分析其对图推理场景的可借鉴性。

\subsection{链式与分解式推理}

Chain-of-Thought 提示通过要求模型先生成中间推理步骤，再给出最终答案，以提升复杂问题求解能力\cite{wei2022chain}。Self-Consistency 在此基础上采样多条思维链并进行投票，降低单次生成的偶然误差\cite{wang2023selfconsistency}。Least-to-Most Prompting 则强调先求解简单子问题，再逐步组合为完整解答\cite{zhou2023least}。Zero-Shot CoT\cite{kojima2022zeroshot}与 Selection-Inference\cite{creswell2023selection}的结果表明，即便在缺少显式示例或需要保持中间步骤可解释的情况下，合理的推理分解也能改善模型表现。这些方法共同说明，对于复杂任务而言，分步思考是提升模型稳定性的有效手段。

然而，这些方法最初主要面向自然语言问答、数学计算或一般数值推理任务。在图推理场景中，模型除了给出中间分析步骤，还必须把这些步骤落实到节点、边、路径、连通性或组合约束等具体图操作上。如果中间过程只有自然语言解释，而缺少算法级的规划表示，模型生成的思路可能仍然比较笼统，后续程序实现也难以准确判断应调用哪类图算法、如何组织输入以及如何检查输出约束。

\subsection{搜索式与结构化推理}

Tree-of-Thoughts 将问题求解过程看作一种搜索过程，通过扩展、评价和回溯多条候选思路，增强模型在复杂任务中的探索能力\cite{yao2023tree}。Graph-of-Thoughts 进一步把不同思路组织成图结构，使中间状态能够被复用，并支持思路之间关系的显式建模\cite{besta2024graph}。ReAct\cite{yao2023react}将推理步骤与外部行动放在同一循环中交替执行；Self-Refine\cite{madaan2023selfrefine}、CRITIC\cite{critic2024}以及过程验证方法\cite{lightman2023processreward}则说明，中间结果生成之后仍可以继续接受检查、反馈和修正。这些研究共同表明，大语言模型的推理过程并不一定要表现为单条线性文本链路，也可以被组织为可搜索、可检查和可修复的结构化过程。

这一点对图推理任务具有直接启发。图任务本身由节点、边及其约束关系构成，天然具有结构化特征。如果中间规划也采用更明确的结构化表达，就可以更自然地连接任务解析、图算法选择、外部知识检索和代码实现等后续环节，减少各阶段之间的语义偏移。

\subsection{推理增强方法的阶段性演进}

若从技术路线演进来看，大语言模型推理增强方法大致可以分为三个阶段。第一阶段主要依靠提示工程，通过示例构造、角色设定和思维链提示等方式，引导模型给出更完整的中间推理过程。第二阶段开始引入外部程序、工具或 API，代表性方法包括 Toolformer\cite{schick2023toolformer}、ToolLLM\cite{qin2023toolllm}和 Gorilla\cite{gorilla2024}，其目标是将精确计算、查询和部分执行操作交给外部环境完成。第三阶段进一步关注多智能体协作和过程控制。CAMEL\cite{camel2023}、MetaGPT\cite{metagpt2023}、AgentVerse\cite{agentverse2023}、AutoGen\cite{wu2023autogen}、Language Agents as Optimizable Graphs\cite{zhuge2024languageagents}、Generative Agents\cite{park2023generativeagents}和 Voyager\cite{wang2023voyager}等工作体现了从单模型回答向多角色协同、过程可控和持续执行的转变；Promptbreeder\cite{zhou2023promptbreeder}则进一步探索了提示策略的自动演化。

图推理任务与上述演进路线具有较高契合度。一方面，图任务具有结构化输入、明确算法定义和可执行求解过程；另一方面，任务结果通常可以通过程序运行或约束检查进行验证。因此，与其要求单个模型在一次生成中同时完成问题理解、算法选择和程序实现，不如将这些功能拆分为多个可检查阶段，并通过中间表示和外部工具提高整体稳定性。对于本文研究的问题，这一路径也为后续框架设计提供了直接依据。

\begin{figure}[htbp]
\centering
\includegraphics[width=0.94\textwidth]{figures/paper/reasoning_evolution.pdf}
\caption{大语言模型推理增强方法的阶段性演进}
\label{fig:reasoning_evolution}
\end{figure}

图~\ref{fig:reasoning_evolution} 概括展示了大语言模型推理增强方法的技术演化过程。对本文而言，这一演化过程不仅是相关研究背景，也说明了图推理框架设计中需要关注的重点：仅通过提示词提升单次回答质量并不充分，还需要将外部知识、执行环境和验证反馈纳入同一求解流程中，从而形成更加稳定的程序化图推理系统。

\section{工具增强与多智能体协同}

由于图推理任务具有可执行、可验证的显著特征，工具调用、程序生成与多智能体协同成为连接语言理解与算法求解的重要桥梁。本节重点讨论相关技术在图任务中的适配方式。

\subsection{程序生成驱动的图任务求解}

PAL 的结果表明，在不少数学推理和符号推理任务中，让模型先生成程序，再由解释器执行计算，通常比直接输出答案更稳定\cite{gao2023pal}。在通用程序生成方向，AlphaCode 展示了大模型处理复杂编程问题的能力\cite{liang2022alphacode}。与此类似，GCoder\cite{zhang2025gcoder}等面向图任务的研究也尝试把自然语言图问题转换为程序，并通过 NetworkX\cite{hagberg2008networkx}等成熟图算法库完成具体求解。

这种程序生成驱动的图任务求解方式，关键在于把精确计算环节从语言模型本身转移到外部执行环境中。语言模型主要负责理解题意、选择算法并组织程序逻辑，具体的图计算、路径搜索或组合优化过程则交由程序和图算法库完成。对于图推理这类既强调结构约束又要求精确结果的任务，这一路线能够减少模型直接生成答案时的随机性和计算误差。

\subsection{工具调用与外部知识使用}

Toolformer\cite{schick2023toolformer}和 ToolLLM\cite{qin2023toolllm}强调模型可以学习何时调用外部工具、使用什么接口以及如何整合工具返回结果。Gorilla\cite{gorilla2024}与 CRITIC\cite{critic2024}的结果表明，模型不仅需要掌握调用能力，还需要在工具交互过程中完成 API 对齐、结果校验和错误纠正。在图推理场景下，外部工具不仅包括图计算函数，还包括图算法文档、函数说明和示例代码等知识源。

NetworkX 是 Python 生态中使用广泛的图计算库，支持丰富的图结构表示与算法接口，包括路径搜索、中心性分析、匹配、生成树、团搜索和流计算等\cite{hagberg2008networkx}。对于语言模型生成程序而言，NetworkX 的 API 语义相对清晰，文档覆盖较完整，适合作为外部知识库的重要来源；同时，它覆盖大量经典图算法，便于将不同基准测试集上的求解过程统一到同一工具生态内。本文将文档知识组织为外部知识库，每条知识包含函数名、适用任务、输入参数、返回形式及示例说明。该设计参考了 DPR\cite{karpukhin2020dpr}、RAG\cite{lewis2020rag}、RETRO\cite{borgeaud2022retro}、Atlas\cite{izacard2023atlas}、Faiss\cite{johnson2019faiss}和 Self-RAG\cite{asai2024selfrag}等检索增强方法，相关技术脉络可参见 RAG 综述\cite{gao2024ragsurvey}。知识库的作用并非直接替代模型推理，而是在模型形成规划后，为代码生成提供可靠的实现参照。

但能够接入工具并不意味着模型一定能够正确使用工具。若模型没有准确判断当前任务对应的算法意图，或者对 API 的功能、参数和返回结果理解不充分，生成的程序仍可能出现函数选择错误、参数传递不当或约束条件遗漏等问题。因此，在工具调用之前加入规划表示，可以先明确求解目标、算法步骤和关键约束，再据此选择和使用外部工具，从而提高工具调用的准确性。

\subsection{图任务中程序生成的适配性分析}

程序生成之所以在图推理任务上具有明显潜力，是因为很多图问题本质上都存在标准化算法流程。例如最短路径问题对应路径搜索算法，最大流问题对应流网络算法，匹配问题对应二分图匹配算法，而团搜索和旅行商问题则可以通过枚举、搜索或优化方法求解。一旦系统能够正确识别任务类型，就有可能借助成熟算法库获得高精度结果。

但图任务中的程序生成也面临比一般代码生成更严格的要求。首先，图构建方式必须正确，错误的方向、边权或时间属性会直接破坏后续求解。其次，许多基准测试集对输出格式极为严格，哪怕节点顺序、集合格式或空格细节不一致，也可能被判为错误。最后，组合优化任务不仅要求程序可运行，还要求解满足可行性和最优性双重约束。近期的 GraphThought\cite{graphthought2025}和伪代码注入方法\cite{pseudocode2025}也尝试缩小语言理解与图算法执行之间的落差。因此，图任务中的程序生成更适合与规划和验证机制结合，而不是依赖一次性生成。

\subsection{多智能体协同求解}

多智能体系统通常将复杂任务拆分给多个功能角色，例如规划者、执行者、检索者、审阅者和修复者等。不同角色在提示词、上下文和目标函数上各有侧重，可以缓解单一智能体既要理解问题、又要生成程序、还要验证结果的压力。GraphTeam\cite{li2024graphteam}面向图分析设置多智能体角色；CAMEL\cite{camel2023}、MetaGPT\cite{metagpt2023}、AgentVerse\cite{agentverse2023}、Generative Agents\cite{park2023generativeagents}与 Voyager\cite{wang2023voyager}则分别从角色扮演协作、软件式协同框架、群体模拟和持续探索等角度展示了多智能体系统的潜力。

不过，现有多智能体方法仍有进一步改进空间。首先，部分方法虽然设置了不同角色，但协作过程更多表现为多轮讨论或意见汇总，缺少一个能够贯穿任务解析、检索、编码和修复阶段的统一规划表示。其次，知识检索和执行验证在一些系统中仍属于辅助环节，通常用于补充信息或检查结果，并没有与结构化求解流程深度绑定。对于图推理任务而言，角色数量的增加并不能直接保证结果可靠；更关键的是让各角色围绕同一规划目标工作，并通过执行反馈形成可回退、可修正的闭环。因此，本文进一步引入规划先行与验证闭环相结合的组织方式，用于提升多智能体图推理过程的一致性和稳定性。

结合已有研究，面向图推理的多智能体系统在设计上至少需要考虑四个方面。第一，角色边界应当明确。问题解析、规划生成、代码实现和执行验证如果职责混杂，系统虽然形式上完成了角色划分，但在出错时仍难以判断问题来自哪个阶段。第二，中间表示需要保持稳定。不同智能体之间不宜完全依赖自由文本对话传递信息，而应通过规划、检索结果、代码和测试反馈等结构化阶段结果完成协作。第三，系统需要为外部知识库、执行环境和图算法库预留统一接口，使后续替换或扩展相关模块时，不必重新设计整体流程。第四，闭环验证应成为系统的核心环节。多智能体协作的目标不应只是生成更多候选答案，而应是在生成后尽快发现错误、定位错误，并根据反馈完成修正。

这几项原则本质上体现了从语言生成系统向可控推理系统的转变。对于图推理任务而言，这种转变尤其重要，因为很多失败并非来自语义理解完全错误，而是来自中间流程缺少稳定约束，导致错误 API 被调用、关键条件被遗漏或格式结果不一致。从功能分工看，问题解析、规划组织、知识检索、程序实现与结果验证可以分别由不同角色承担。这种按角色划分局部目标的方式，有助于降低一次性端到端生成的风险，也为后续章节讨论规划反馈闭环提供了协作基础。

\section{相关技术总结}

前述各节分别从图推理任务基础、大语言模型推理增强、工具增强与多智能体协同等角度梳理了本文工作的技术背景。为了明确相关工作的整体位置和本文方法的切入点，本节从方法分类、任务能力分解和相关工作比较三个层面对已有技术进行归纳。

\subsection{图推理相关方法分类}

公开基准测试集在图推理研究中扮演了重要角色。一方面，它们将图算法问答、指令执行和动态图查询等任务组织为可重复比较的标准测试集，使研究者能够更客观地评价模型性能。另一方面，不同基准测试集各有侧重，推动了图推理研究从单一任务优化转向跨任务泛化能力评估。GraphArena\cite{tang2025grapharena}强调经典图算法与组合优化任务，NLGraph\cite{chen2024nlgraph}关注自然语言描述下的图推理能力，GraphWiz\cite{gong2024graphwiz}和 GraphInstruct\cite{graphinstruct2024}侧重指令式图任务求解，LLM4DyG\cite{llm4dyg2024}则进一步引入动态图及时间演化场景。

在本文中，引入基准测试集的目的并不只是说明使用了哪些数据，更重要的是分析不同基准分别考察模型的哪些能力。各数据集在输入形式、图结构复杂度、约束强度和输出格式要求上并不相同。只有先区分这些差异，才能解释同一种方法为什么可能在某个基准上取得较好结果，却在另一个基准上出现性能下降。因此，本文在后续实验中不仅报告总体结果，也进一步从任务类别、样本难度和扰动鲁棒性等角度进行分析。

为了更清楚地展示相关方法之间的关系，图~\ref{fig:llm_graph_reasoning_taxonomy} 从两个维度对代表性工作进行了整理。横轴表示方法所面向的任务属性，左侧更接近通用推理任务，右侧更接近专门的图推理任务；纵轴表示方法的系统组织程度，下方主要对应单次生成或弱反馈方法，上方则对应系统化协同或强闭环方法。位于左下区域的 CoT\cite{wei2022chain}、Self-Consistency\cite{wang2023selfconsistency}、ToT\cite{yao2023tree}、GoT\cite{besta2024graph}和 CRITIC\cite{critic2024}等方法，主要属于通用推理增强范式，关注如何通过思维链扩展、路径搜索或自我修正提升单模型推理能力。位于右下区域的 NLGraph\cite{chen2024nlgraph}、GraphWiz\cite{gong2024graphwiz}、GraphGPT\cite{tang2024graphgpt}、GraphInstruct\cite{graphinstruct2024}、GraphLLM\cite{chai2023graphllm}、GCoder\cite{zhang2025gcoder}和 LLM4DyG\cite{llm4dyg2024}等工作，则更集中于图任务本身，重点讨论图结构表示、图算法调用、程序化求解以及动态图推理等问题。

\begin{figure}[htbp]
\centering
\includegraphics[width=0.96\textwidth]{figures/paper/llm_graph_reasoning_taxonomy.pdf}
\caption{图推理相关方法分类示意图}
\label{fig:llm_graph_reasoning_taxonomy}
\end{figure}

右上区域的方法通常更重视工具使用、多智能体协作以及闭环控制。工具增强方法包括 Toolformer\cite{schick2023toolformer}、ToolLLM\cite{qin2023toolllm}和 Gorilla\cite{gorilla2024}。多智能体方法包括 GraphTeam\cite{li2024graphteam}、GraphAgent-Reasoner\cite{hu2024graphagentreasoner}、CAMEL\cite{camel2023}、MetaGPT\cite{metagpt2023}、AgentVerse\cite{agentverse2023}和 MA-GTS\cite{yuan2025magts}。这些研究已经不再只关注单次回答是否更准确，而是开始关注任务求解过程能否被组织、回退和协同执行。此外，DPTS\cite{ding2025dpts}、HyperTree Planning\cite{gui2025hypertree}、Graph-enhanced Plan Reasoning\cite{lin2024plangraph}、ARise\cite{zhang2025arise}、Decision Pivots\cite{cho2025decisionpivots}和 ReAct\cite{yao2023react}等工作位于图中上半区域，说明越来越多方法开始通过显式搜索、层次化规划、图式计划表示和风险感知推理来提高复杂任务求解过程的稳定性。总体来看，该二维分类图展示了两个相互交织的发展方向：一是从通用推理增强逐步转向图任务适配，二是从单模型单次生成逐步转向系统化协同与闭环修正。本文方法同时关注图任务适配和系统化闭环求解，因此与这两条发展方向保持一致。
\subsection{图推理任务分类与能力分解}

为了更系统地理解图推理问题的研究范围，本文将常见任务划分为五类：结构查询类、路径搜索类、排序与可达性类、网络流与匹配类以及组合优化类。结构查询类任务通常关注邻居、度数、连通分量和局部关系判断，更依赖图构建的准确性；路径搜索类任务关注最短路径、可达性与遍历顺序，要求系统稳定处理边权、方向和路径格式；排序与可达性类任务强调全局先后约束，例如拓扑排序和有向无环图判定；网络流与匹配类任务要求对容量、配对约束和全局最优性同时建模；组合优化类任务如最大团、最小点覆盖和旅行商问题，则提高了搜索复杂度和验证难度。

\begin{table}[htbp]
\centering
\small
\caption{图推理任务分类及其核心能力要求}
\label{tab:task_taxonomy}
\renewcommand{\arraystretch}{1.15}
\begin{tabularx}{\textwidth}{
>{\raggedright\arraybackslash}p{2.9cm}
>{\raggedright\arraybackslash}X
>{\raggedright\arraybackslash}X}
\toprule
任务类别 & 代表任务 & 核心能力要求 \\
\midrule
结构查询类 & 邻居查询、度数计算、连通分量 & 正确构图、读取局部关系、规范输出 \\
路径搜索类 & 最短路径、DFS/BFS、可达性判断 & 处理边权与方向，生成合法路径 \\
排序与可达性类 & 拓扑排序、环检测、前驱查询 & 建模全局依赖，满足顺序约束 \\
网络流与匹配类 & 最大流、二分图匹配、最小割 & 处理容量、匹配与全局最优约束 \\
组合优化类 & 最大团、最小点覆盖、旅行商问题 & 满足可行性约束，并验证最优性 \\
\bottomrule
\end{tabularx}
\end{table}

表\ref{tab:task_taxonomy}所呈现的任务分类说明，图推理并不是一个单一难点，而是由图表示、约束识别、算法匹配、程序实现、结果验证多个能力共同决定的复合问题。对大语言模型而言，不同任务类别暴露出的短板也并不相同。例如，结构查询类任务更容易受到输入解析错误影响，而组合优化类任务更容易在搜索空间控制和最优性验证阶段失效。正因为这种差异性存在，本文在方法设计上没有将所有图任务统一视为自然语言问答，而是选择通过规划与验证机制显式刻画任务间的结构差异。

\subsection{相关工作对比}

为更清晰地总结现有图推理方法之间的差异，表\ref{tab:related_compare}从推理方式、是否使用外部知识、是否执行程序以及是否具备验证闭环四个维度进行了对比。表中所涉及的代表性技术路线分别对应直接提示推理、结构感知提示、程序生成和多智能体协作等研究方向，其代表工作包括 ReAct\cite{yao2023react}、GraphGPT\cite{tang2024graphgpt}、GCoder\cite{zhang2025gcoder}和 GraphTeam\cite{li2024graphteam}。

\begin{table}[htbp]
\centering
\small
\caption{相关图推理方法对比}
\label{tab:related_compare}
\renewcommand{\arraystretch}{1.15}
\begin{tabularx}{\textwidth}{
>{\raggedright\arraybackslash}p{2.4cm}
>{\raggedright\arraybackslash}X
>{\centering\arraybackslash}p{1.8cm}
>{\centering\arraybackslash}p{1.8cm}
>{\centering\arraybackslash}p{2.1cm}}
\toprule
方法类别 & 代表思路 & 外部知识 & 程序执行 & 验证闭环 \\
\midrule
直接提示推理 & 通用推理增强 & 否 & 否 & 否 \\
结构感知提示 & 图结构提示或图任务对齐方法 & 可选 & 否 & 弱 \\
程序生成方法 & 将图任务转化为代码求解 & 可选 & 是 & 弱 \\
多智能体方法 & 多角色协作完成规划、生成与修复 & 可选 & 可选 & 部分具备 \\
本文方法 & 规划引导检索、执行与反馈修复 & 是 & 是 & 强闭环 \\
\bottomrule
\end{tabularx}
\end{table}

表中比较的是方法机制层面的典型特征，而非对所有具体实现作绝对判定。其中，外部知识指是否显式引入文档、样例库、代码库或检索知识；程序执行指候选程序是否被真实运行；验证闭环指系统能否依据执行结果自动触发修复、补充检索或上游回退。若方法仅进行单次程序执行，或依赖角色讨论完成检查但缺少稳定的自动反馈路径，本文将其归为弱反馈或回退有限。

由表\ref{tab:related_compare}可见，现有方法之间的差异并不只体现在是否使用图结构这一单一维度上，而更多体现在是否引入外部知识、是否真正执行程序，以及是否形成可回退的验证闭环。直接提示推理与结构感知提示方法更关注输入表达和单次生成质量，程序生成方法开始借助外部执行提升精确性，多智能体方法则尝试通过角色协作增强复杂任务处理能力，但在统一规划表示、检索支撑和反馈推理的协同组织上仍有整合空间。基于这一观察，本文方法进一步围绕显式规划、知识检索、程序执行与反馈推理构建协同机制。

\section{本章小结}

本章从图推理任务基础、大语言模型推理增强、工具增强与多智能体协同以及相关技术总结等方面，梳理了本文工作的理论基础。综合现有研究可见，图推理问题同时涉及结构解析、算法选择、程序执行和结果校验，单次文本生成难以稳定覆盖这些环节。因此，后续方法需要进一步围绕显式规划、知识检索和验证闭环组织求解过程。下一章将对所提方法进行详细说明。
